\documentclass{article}
\usepackage{graphicx} % Required for inserting images
\usepackage{amsmath}
\usepackage{amsfonts}
\usepackage{eufrak}
\usepackage{hyperref}
\title{Scalar Field HDN}
\author{Alfred Ricker}
\date{November 2025}

\newcommand{\D}{\mathrm{d}}

\begin{document}

\maketitle

\section{Overview}
The essential properties of the Scalar Field Hyperbolic Dynamic Network (HDN) that I have outlined in my notes are as follows.
\begin{itemize}
    \item There is a series of ``input'' nodes that receive input information. As of now they are just simple unsigned integers that encode characters into bits.
    \item The nodes exist in a Hyperbolic space $\mathbb{D}$
    \item The nodes have a binary firing function, that when activated ``release'' energy so that $v(t_f) = 0$
    \item Firing triggers an increase in all other node energies proportional to the interaction kernel $K(d)$ and \textit{active} neurons move in the direction of the fired node (with $dz \propto K(d)$ as well). 
    \item Reinforcement mechanism to reward good outputs and punish bad ones.
\end{itemize}

\subsection{Modes of Thought}
In the human brain, there seems to be unconscious activity and conscious activity, and these two modes guide each other. For example, I can think about what I am going to say, or sometimes I automatically respond to questions without giving it much thought. There is a layered mechanism of decision making occurring. Does this input signal require thought? Do my thoughts align with a sensible output? Do I output now or give it more thought?

It seems that the inputs that map to some abstract understanding map to some analogous space of outputs.

\subsection{Understanding Instructions and Memory}
The network should attempt to figure things out that it doesn't understand. It seems that we need to somehow ``remember'' activity in regions close to the input nodes so that these may be re-excited when a given input does not produce a confident output or equilibrium state within the system. \\

As for a system of memory, this does occur somewhat naturally within the mechanics of energy, gravity, and pathways of decreasing plasticity (as will be explained in future sections). An input that is seen many times is essentially remembered by activating the appropriate pathway and reaching a steadfast equilibrium state. However, in the brain the recollection mechanism seems to be quite complicated, as focused pathways seem to activate meaningful thoughts.

\section{Algebraic Motivation}
Here I will overview the motivations for the algebra of the neural network more in depth.
\subsection{Hyperbolic Space}
Hyperbolic space allows for a hierarchical representation of concepts. The shortest distance between two nodes in this space, called a geodesic in the language of differential geometry, curves inward towards the origin (it is orthogonal to the boundary $\partial \mathbb{D}$). Since the size of space increases exponentially as we move away from the origin, more specific concepts will be pulled towards the nearest less specific concepts that live closer to the origin. This is a natural result of the increasing size of space, it is significantly less probable for two complex thoughts to feel the energy of one another.

\subsection{Neuron Variables}
A neuron is represented by the following variables
\begin{itemize}
    \item $z$: the position in the space. The space has a well defined distance function between points $d(z_1,z_2)$
    \item $\omega$: The frequency of firing.
    \item $v$: The internal energy of the neuron
    \item $\rho$: The energy threshold of firing.
    \item $\tau$: The time elapsed since firing.
    \item $\eta$: Plasticity or coefficient of freedom to move throughout the space.
\end{itemize}
Therefore the set of neurons is defined as 
\begin{equation}
    \mathcal{N} = \{ (z_1,\omega_1, \nu_1), ... ,(z_m,\omega_m, \nu_m) \}
\end{equation}
The space has negative curvature and a radial symmetry, i.e. neurons located at $|z| = c$ are all the same distance from the origin but distances grow larger for the same angles as we move further away from the origin. This allows for a natural mapping of $|z|$ to frequency and coincides with Zipf's law, where we can exhibit a tree like structure as neurons move further from the origin. 

\subsection{Global Variables}
\begin{itemize}
    \item $T$: The total time elapsed since the start of the simulation.
    \item $E$: Total energy of the space $\mathbb{D}$
\end{itemize}
Global variables are important for deciding things such as how to create new neurons, when to decide convergence of thought on a particular problem, etc.

\subsection{Dimensionality}
It is worth examining the tradeoffs of whether to make the space larger with added dimensions to allow for more linear independence among certain objects. The distance formula between two points $x,y$ in any dimensional hyperbolic space is
\begin{equation} 
d_{\mathbb{D}}(x, y) = \mathrm{arccosh} \left( 1 + 2 \frac{||x - y||^2}{(1 - ||x||^2)(1 - ||y||^2)} \right) 
\end{equation}
Let $c$ be a central concept (a point in the manifold) and let $b_i$ be boundary concepts. How many $b_i$ can exist such that $d(c,b_i) < \varepsilon \, \forall \, b_i $ and $d(b_i,b_j) > \gamma \, \forall \, (b_i,b_j)$. Let's start by restricting $d(c,b_i) = \varepsilon$. The equation for the area of a hypersurface on a manifold is 
\begin{equation}
    A(\Sigma) = \int_\mathcal{U} \sqrt{\det(h)} \, \D u_1 \D u_2...\D u_n
\end{equation}
Where $h$ is the induced metric tensor on the surface
\begin{equation}
    h_{ij} = \sum_{\mu, \nu} \frac{\partial x^\mu}{\partial u^i} \frac{\partial x^\nu}{\partial u^j} g_{\mu \nu} 
\end{equation}


\section{Describing the Variables}
\subsection{Interaction Kernel}
For hyperbolic space, the standard interaction from Euclidean geometry of $\frac{1}{r^2}$ fails to apply for the exponential expansion in space with respect to $r$. If we apply the solution to the heat equation in $\mathbb{H}$, we get an interaction kernel
\begin{equation} 
K(d) \propto e^{-\frac{d_{\mathbb{D}}^2}{4t}} \cdot \left( \frac{d_{\mathbb{D}}}{\sinh(d_{\mathbb{D}})} \right)^{\frac{n-1}{2}} 
\end{equation}
This kernel is computationally expensive, so a preferable kernel is a simpler one that maintains stability. Let the effect that a neuron $n_i$ located at $x$ have on an arbitrary point $z$ be defined by
\begin{equation} K(d) = e^{-\beta d_{\mathbb{D}}(x, z)} \end{equation}
Where $\beta > D-1$ to avoid integral divergences. I will set $\beta = D$. \\

It is also a desirable feature to have interactions between neurons be dependent on their ``activity'' as to introduce a factor that decays with the time since firing $\tau$.
\begin{equation}
    K(d,\tau) =e^{-\beta d(x,z) - \alpha\tau}
\end{equation}
We will call this kernel the tau dependent kernel.

\subsection{Internal Energy}
Each neuron has an internal energy that increases as local neurons fire and decays as there is low local energy (firing activity). We denote this internal energy of neuron $n_j$ as $v_j$. The most important feature of the internal energy $v_j$ is that it is excited when neuron $n_l$ fires by a factor that is proportional to the interaction kernel. We will not use the tau dependent interaction kernel, as we want neurons to receive signals independent of their time since firing.
\begin{equation}
    v'_j \propto K(d_{jl} )
\end{equation}

The internal energy should have a buffer period so that it is hard for a neuron to build up energy quickly after firing. This prevents a local region from firing in loop, collapsing the whole system into a diverging energy black hole. Let $t_f$ be the time that the neuron fired last. We introduce a factor
\begin{equation}
    \xi = e^{-b/(1+\tau)^2}
\end{equation}
Therefore directly after firing, we decrease the energy gained by a factor of $e^{-b}$ and this damping factor goes away quadratically with time since firing $t-t_f$. If the neuron has not fired at all, we set $t-t_f \to \infty \implies \xi = 1$. The equation for internal energy updates is 
\begin{equation}
    \frac{\D v_j}{\D t} = e^{-b/(1-\tau)^2}\sum_{l \in \mathrm{Firing}} e^{-\beta d_{jl}}
\end{equation}

The firing condition of a neuron is $v \ge \rho$. This resets the energy $v' = v-\rho$ where the neuron can be thought of as releasing its energy to the system.



\subsection{Manifold Energy and Neurogenesis}
To determine locations for concept formation (neurogenesis), we distinguish between \textit{Field Density} and \textit{Semantic Consensus}.

The \textbf{Field Density} $E(z)$ measures the raw activity at a point:
\begin{equation}
    E(z) = \sum_{n_i \in \mathcal{N}} K(d(z,n_i), \tau_i)
\end{equation}
This function spikes at existing neuron locations. To identify the "center" of a distributed thought pattern (where a new concept \textit{should} exist but perhaps does not yet), we utilize the \textbf{Weighted Fréchet Energy} $V(z)$:
\begin{equation}
    V(z) = \sum_{n_i \in \text{Active}} w_i d_{\mathbb{D}}^2(z, n_i)
\end{equation}
The minimizer of this function, $\mu = \arg\min_z V(z)$, represents the \textit{Hyperbolic Barycenter} of the active concepts. 

\textbf{Creation Mechanism:} It is desirable to have some mechanism of Homeostatic Plasticity / Neurogenesis so the network may expand the more it learns. For the network, an unfamiliar concept can be measured by the rate of energy decay for firing in the input space. Let's say there is some $E_{|z| < \kappa}$. If 
\begin{equation}
    \frac{dE_{|z|>\kappa}}{dt} < \Xi
\end{equation}
Then that means the rate of energy dispersion beyond the boundary $\kappa$ is lower than our constant $\Xi$ and the message didn't activate significant networks of energy. This should precipitate the creation of a new neuron at $\arg\min_z V(z)$. \\

\textbf{Local Equilibrium}: It is also desirable to have the network have a defined equilibrium within some region $R$ to activate actions (such as triggering memory regions, outputs, etc). 
\begin{equation}
    \int_R \text{Var}(E)dt < \varsigma
\end{equation}


\subsection{Plasticity}
The plasticity $\eta$ of a neuron is a factor that influences how difficult it is for the neuron to move in any direction of the space. This is quite useful because we don't want the network to be exceedingly dynamic. Certain regions and trees of the network will belong to some well learned concepts and actions, which we don't want moving around too much with new activity. Therefore, we can decrease the plasticity with time that a neuron has existed, and the number of times it has fired.
\begin{equation}
    \eta = Ce^{-cT - \mu \omega}
\end{equation}

\subsection{Drift}
There should be some centripetal force acting on the neurons in the system, slowly pushing them out towards the boundary as they are not used (while decreasing $\rho$ in the process because energy signals become weaker at further distances from the origin). When a neuron fires, it should feel a pull towards the origin, acting against the centripetal force. This helps to naturally form the frequency hierarchy and allow for commonly fired nodes to be in proximity to input neurons. \\

The drift parameter of course should depend on plasticity $\eta$ and distance from the origin $|z|$. 

\begin{equation}
    \mathfrak{d} = \beta\eta - \alpha e^{-\mu\tau}
\end{equation}
When $\tau$ is small this drift vector points towards the origin, otherwise it drifts of quadratically dependent on the plasticity (remember that a smaller plasticity means that it is harder for the neuron to move). The drift is yet another component of the differential equation of motion
\begin{equation}
    \frac{dz}{dt} = \lambda \eta(\mathfrak{d} + V(z,t))  
\end{equation}

\subsection{Input Neurons}
Input neurons are to be randomly initialized before the simulation starts within some subspace $|z| < \kappa$, with each base unit of information that can be passed to the machine being mapped to $q$ input neurons randomly distributed within the subspace. This allows for more complex dependencies on the inputs within the network. \\

Input neurons do not have internal energies. They simply receive inputs and fire, releasing energy and dispersing the message to the system.

\section{Actions and Reinforcement}

\subsection{Action Sectors}
Unlike the discrete input nodes, the output of the network is determined by the accumulation of energy in designated \textit{Action Sectors} near the boundary of the manifold, denoted $\mathcal{A} = \{A_1, \dots, A_k\}$. The energy of an action sector is the integral of field density over its region:
\begin{equation}
    E_{A_k}(t) = \int_{z \in A_k} E(z, t) \, \D z
\end{equation}
The network selects an action probabilistically based on a Softmax distribution scaled by a temperature parameter $\Theta$:
\begin{equation}
    P(A_k) = \frac{\exp(E_{A_k} / \Theta)}{\sum_j \exp(E_{A_j} / \Theta)}
\end{equation}

\subsection{Dopaminergic Modulation}
The system learns via a global scalar reward signal $R(t) \in [-1, 1]$. This signal does not directly create energy but modulates the structural variables of neurons that were active during the decision window.

\subsubsection{Plasticity Annealing}
To enforce learning, we utilize a mechanism where reward ``freezes'' the network configuration, while punishment introduces volatility (simulating a search for a lower energy state). The plasticity $\eta$ updates as:
\begin{equation}
    \frac{\D \eta}{\D t} = -\gamma_{\eta} \cdot R(t) \cdot \eta
\end{equation}
A positive reward decreases plasticity (crystallization), while a negative reward increases plasticity, allowing neurons to drift out of erroneous local minima.

\subsubsection{Threshold Adaptation}
We implement a Hebbian-like rule for the energy threshold $\rho$. Active nodes that lead to reward become more sensitive:
\begin{equation}
    \rho_{new} = \rho_{old} - \lambda R(t) \frac{E_{local}(z)}{E_{max}}
\end{equation}
This ensures that pathways leading to successful predictions are easier to traverse in future iterations.

\subsubsection{Automated vs. Manual Modulation}
We can develop a train set that has the network perform actions and assigns a dopamine signal to its outputs, analogous to the backpropagation signals of standard ANNs. A more advanced system would be to train the network to be able to interpret language feedback to what it is outputting, as some sort of social reward function.

\end{document}
